% E4: Does the flow-time conditioning do work, and what does it cost?
% Drop-in section. Uses only what the main paper already loads.
% All values are filled from discovery run counterfactual-20260924-154113.
% Artefact keys for each table are named in a comment above it.

\providecommand{\RESULT}[1]{\textbf{[#1]}}

\section{What the Flow-Time Conditioning Buys, and What It Costs}
\label{sec:tau}

Conditioning each transcoder on the flow timestep $\tau$ is the architectural
choice that distinguishes these sparse codes from an ordinary transcoder, and
it has a price that is easy to miss. A feature is identified by the triple
$(\ell,\tau,i)$, so the same dictionary index at ten denoise steps is ten
computational nodes, and a circuit traced over them is traced over a node space
ten times the dictionary. Neither side of that trade has been measured. This
section measures both from a feature-discovery run, with no retraining and no
additional inference.

The finding is that the trade is real but uneven. Conditioning does genuine
work in the middle of the action expert and at its last layer, and almost none
at the first two layers or at layers twelve through sixteen. Meanwhile the cost
is uniform: at every one of the eighteen layers, codes at adjacent flow times
are near-duplicates.

\subsection{Hypothesis}

\paragraph{H4.} The sparse code depends on the flow timestep in a way that
smooth interpolation between denoise steps does not explain.

\emph{Falsified if} the code at distant flow times is as correlated as a smooth
drift along $\tau$ would predict, which would mean the time MLP reparameterises
the code without restructuring it.

\emph{Decision quantity.} For layer $\ell$, let $\mu_{\ell,\tau}\in\mathbb{R}^{d_{\mathrm{features}}}$
be the per-feature mean activation at flow time $\tau$, and let
$r_{\ell}(\tau_a,\tau_b)$ be the correlation across features between
$\mu_{\ell,\tau_a}$ and $\mu_{\ell,\tau_b}$. Write $r^{\mathrm{adj}}_{\ell}$ for
the mean over adjacent pairs and $r^{\mathrm{far}}_{\ell}$ for the first-to-last
pair. The decision quantity is the \emph{structure ratio}
\begin{equation}
\rho_{\ell}
=
\frac{r^{\mathrm{far}}_{\ell}}
     {\left(r^{\mathrm{adj}}_{\ell}\right)^{\,n-1}},
\label{eq:tau-structure}
\end{equation}
with $n$ the number of denoise steps.

\emph{Rule, fixed in advance.} A layer \emph{uses its conditioning} when
$\rho_{\ell} < 0.5$ and the layer is reliable in the sense of
Section~\ref{sec:tau-method}. Verdict set: \{conditioning used at every usable
layer, used at some, unused, not measured\}.

\subsection{Method}
\label{sec:tau-method}

\paragraph{The null is smooth drift, not zero correlation.}
A low distant-$\tau$ correlation on its own is not evidence of structure. If the
code simply wanders a little at each denoise step, and each step's change is
independent of the last, then correlation decays geometrically with lag and the
first-to-last pair sits at $(r^{\mathrm{adj}})^{n-1}$. That is what a time MLP
doing nothing but smoothly reparameterising would produce, and at
$r^{\mathrm{adj}}=0.98$ over ten steps it already predicts a distant correlation
of $0.83$. The denominator of Eq.~\ref{eq:tau-structure} is that prediction, so
$\rho$ asks whether the code changes with $\tau$ by more than drift, which is
the question the architecture's claim actually turns on.

\paragraph{The reliability ceiling.}
These are correlations between means estimated from finitely many observations,
so sampling noise attenuates every one of them and even two identical flow times
would not correlate at one. With $s_{\ell,\tau,i}$ the per-feature standard
deviation over $N$ observations, the estimated means carry a sampling variance
of $\mathbb{E}_i[s^2]/N$ and the reliability is
\begin{equation}
\kappa_{\ell,\tau}
=
1 - \frac{\mathbb{E}_i\!\left[s_{\ell,\tau,i}^2\right]/N}
         {\mathrm{Var}_i\!\left[\mu_{\ell,\tau,i}\right]},
\label{eq:tau-reliability}
\end{equation}
so a correlation divided by $\sqrt{\kappa_a \kappa_b}$ is the disattenuated
estimate. A layer whose minimum reliability falls below $0.2$ is reported as not
measurable at this sample size rather than corrected, because dividing by a
near-zero ceiling manufactures structure. Features that never fire at any flow
time carry no information about $\tau$ and are excluded, so the correlation is
taken over the layer's live dictionary.

\subsection{Results}

% Filled from tau_conditioning.csv and tau_conditioning.json -> layers, summary.
\begin{table}[t]
\caption{Flow-time structure by layer, from $160$ observations at $n=10$ denoise
steps. Drift is the smooth-interpolation prediction
$(r^{\mathrm{adj}})^{9}$; $\rho$ is Eq.~\ref{eq:tau-structure}. Reliability is
the minimum of Eq.~\ref{eq:tau-reliability} over the layer's flow times.}
\label{tab:tau-layers}
\begin{center}
\small
\begin{tabular}{ccccccc|c}
\hline
Layer & Live & $\kappa_{\min}$ & $r^{\mathrm{adj}}$ & $r^{\mathrm{far}}$ & Drift & $\rho$ & Uses $\tau$ \\
\hline
0  & $3428$  & $0.99$ & $0.989$ & $0.654$ & $0.908$ & $0.72$ & no \\
1  & $4230$  & $0.99$ & $0.989$ & $0.647$ & $0.903$ & $0.72$ & no \\
2  & $4923$  & $0.99$ & $0.979$ & $0.157$ & $0.823$ & $0.19$ & yes \\
3  & $7551$  & $0.99$ & $0.977$ & $0.116$ & $0.811$ & $0.14$ & yes \\
4  & $15291$ & $1.00$ & $0.980$ & $0.239$ & $0.832$ & $0.29$ & yes \\
5  & $15234$ & $1.00$ & $0.980$ & $0.307$ & $0.836$ & $0.37$ & yes \\
6  & $16166$ & $0.99$ & $0.977$ & $0.204$ & $0.814$ & $0.25$ & yes \\
7  & $14989$ & $0.99$ & $0.978$ & $0.248$ & $0.822$ & $0.30$ & yes \\
8  & $16253$ & $0.99$ & $0.976$ & $0.174$ & $0.806$ & $0.22$ & yes \\
9  & $16300$ & $0.98$ & $0.976$ & $0.176$ & $0.804$ & $0.22$ & yes \\
10 & $16158$ & $0.98$ & $0.979$ & $0.326$ & $0.824$ & $0.40$ & yes \\
11 & $16309$ & $0.98$ & $0.977$ & $0.317$ & $0.808$ & $0.39$ & yes \\
12 & $16135$ & $0.98$ & $0.989$ & $0.683$ & $0.908$ & $0.75$ & no \\
13 & $15761$ & $0.98$ & $0.989$ & $0.667$ & $0.908$ & $0.74$ & no \\
14 & $15378$ & $0.98$ & $0.996$ & $0.907$ & $0.962$ & $0.94$ & no \\
15 & $15507$ & $0.98$ & $0.994$ & $0.843$ & $0.947$ & $0.89$ & no \\
16 & $11405$ & $0.99$ & $0.995$ & $0.820$ & $0.952$ & $0.86$ & no \\
17 & $7083$  & $0.99$ & $0.980$ & $0.388$ & $0.834$ & $0.47$ & yes \\
\hline
\multicolumn{8}{l}{$r^{\mathrm{adj}}$: median $0.980$, range $0.976$ to $0.996$. $r^{\mathrm{far}}$: median $0.321$, range $0.116$ to $0.907$.} \\
\multicolumn{8}{l}{Every layer is reliable ($\kappa_{\min}\ge0.98$), so none is excluded and no correction exceeds $2\%$.} \\
\hline
\end{tabular}
\end{center}
\end{table}

\paragraph{H4 verdict: supported at $11$ of $18$ layers.}
Conditioning restructures the code well beyond drift at layers $2$ through $11$
and at layer $17$, where $\rho$ runs from $0.14$ to $0.47$. It does not at
layers $0$, $1$ and $12$ through $16$, where $\rho$ runs from $0.72$ to $0.94$
and the observed distant correlation is close to what smooth interpolation alone
predicts. The split is not a power artefact: reliability is at least $0.98$
everywhere, so the $160$-observation run is ample for this measurement and the
disattenuation changes no correlation by more than two per cent.

\paragraph{The cost is uniform, and it is a circuit-tracing problem.}
At every layer, codes one denoise step apart correlate at $0.976$ or better.
Whatever conditioning does across the full range of $\tau$, it does almost
nothing between neighbours. The consequence is not aesthetic. Because a feature
is a $(\ell,\tau,i)$ triple, two adjacent nodes in a traced circuit can be the
same computational object counted twice, and a backward trace with a budget
measured in nodes spends most of that budget re-expanding near-duplicates. This
is a concrete mechanism for the fan-out reported in
Section~\ref{sec:counterfactual}, where a trace from a single target reached
$1681$ parents. It also suggests a cheap remedy worth testing: collapse or
deduplicate adjacent-$\tau$ nodes before expanding the frontier.

\subsection{Relation to the rest of the paper}

\paragraph{Architecture.} This is the first direct evidence that the time
conditioning changes what the code represents rather than merely how it is
parameterised, and it localises that evidence to the middle of the action
expert. It supports the architectural choice while narrowing the claim: at
layers $14$ to $16$ the conditioning is close to inert, and a cheaper
unconditioned transcoder might serve there.

\paragraph{Feature discovery.} The discovery table lists
$\mathtt{L12}\!:\!\tau\!=\!1.0\!:\!\mathtt{F7584}$ and
$\mathtt{L12}\!:\!\tau\!=\!0.9\!:\!\mathtt{F7584}$ as separate top-ranked
entries with the same top-$20$ context, and reads their co-occurrence as
evidence that the feature is not a single-timestep artefact. The measurement
here gives that reading its footing and its limit: at layer $12$ adjacent flow
times correlate at $0.989$, so the two entries are near-duplicates by
construction. Their agreement confirms stability, but it is close to
automatic and should not be counted as independent support.

\paragraph{Circuit construction.} The redundancy above bears directly on trace
size and on how a node budget should be spent. It does not challenge any
reported circuit; it explains a mechanism by which traces inflate.

\subsection{Limitations}

The measurement is correlational and taken over per-feature means, so it
characterises how the layer's code as a whole moves with $\tau$, not whether any
individual feature keeps its meaning. It uses one discovery run on one dataset
slice; the structure ratio is a summary of the first-to-last pair and does not
describe the shape of the path between them. And it says what conditioning
\emph{does}, not what it is \emph{worth}: a matched transcoder trained with the
time MLP frozen at its zero initialisation would establish whether the
restructuring improves reconstruction at equal sparsity, and remains the
confirmatory experiment.
